\documentclass[11pt,a4paper]{article}

\usepackage[margin=0.8in]{geometry}
\usepackage{graphicx}
\usepackage{booktabs}
\usepackage{tabularx}
\usepackage{listings}
\usepackage{xcolor}
\usepackage{float}
\usepackage{caption}
\usepackage{subcaption}
\usepackage{hyperref}
\usepackage[T1]{fontenc}

\hypersetup{colorlinks=true, linkcolor=blue, urlcolor=blue}

\lstset{
    language=C++,
    basicstyle=\ttfamily\small,
    keywordstyle=\color{blue!60!black}\bfseries,
    commentstyle=\color{gray}\itshape,
    breaklines=true,
    frame=single,
    columns=fullflexible
}

\newcolumntype{Y}{>{\raggedright\arraybackslash}X}
\newcommand{\graphs}[1]{%
  \begin{figure}[H]
    \centering
    \begin{subfigure}[b]{0.32\textwidth}\centering
      \includegraphics[width=\linewidth,height=8cm,keepaspectratio]{trees/#1_cfg.pdf}
      \caption{CFG}
    \end{subfigure}\hfill
    \begin{subfigure}[b]{0.32\textwidth}\centering
      \includegraphics[width=\linewidth,height=8cm,keepaspectratio]{trees/#1_dom.pdf}
      \caption{Dominator tree}
    \end{subfigure}\hfill
    \begin{subfigure}[b]{0.32\textwidth}\centering
      \includegraphics[width=\linewidth,height=8cm,keepaspectratio]{trees/#1_postdom.pdf}
      \caption{Post-dominator tree}
    \end{subfigure}
  \end{figure}}

\title{\textbf{Mini Assignment: SSA, Dominators and Post Dominators}\\
\large Question 1: Dominators and Post Dominators}
\author{CS3423 -- Compilers II}
\date{November 2026}

\begin{document}
\maketitle

\section{Dominators vs.\ post-dominators}

A block $A$ \emph{dominates} block $B$ if you can't reach $B$ from the
entry without passing through $A$ first. A \emph{post-dominator} is the
same idea run backwards: $A$ post-dominates $B$ if every path from $B$
to the exit goes through $A$. So dominators tell you what has definitely
run \emph{before} a block, and post-dominators tell you what will
definitely run \emph{after} it.

\section{How the trees were made}

Each program was compiled to LLVM IR with
\begin{lstlisting}[language=bash]
clang++ -std=c++17 -S -emit-llvm -O0 -Xclang -disable-O0-optnone \
        -fno-discard-value-names N.cpp -o IR/N.ll
\end{lstlisting}
\texttt{-fno-discard-value-names} keeps readable block names like
\texttt{if.then} and \texttt{for.cond} instead of \texttt{\%5}. The trees
were then printed with
\texttt{opt -passes='print<domtree>,print<postdomtree>'}
(output saved in \texttt{trees/N.txt}) and drawn with LLVM's
\texttt{dot-cfg-only}, \texttt{dot-dom-only} and \texttt{dot-post-dom-only}
passes. Running \texttt{./run.sh} redoes all of this.

A few notes for reading the tables:
\begin{itemize}
  \item Every block dominates and post-dominates itself, so the lists only
        show the \emph{strict} relations.
  \item \textbf{idom} / \textbf{ipdom} are the immediate dominator and
        immediate post-dominator, i.e.\ the parent in each tree.
  \item LLVM puts a virtual \texttt{exit} node at the root of the
        post-dominator tree, which is why the function's last block has
        \texttt{exit} as its ipdom.
  \item Only the interesting function in each file is shown. \texttt{main}
        is a single block, so its trees are trivial.
\end{itemize}

%=====================================================================
\newpage
\section{Program 1: if / else-if / else ladder}

\begin{lstlisting}
int classify(int x) {
    int result;
    if (x < 0) {
        result = -1;
    } else if (x == 0) {
        result = 0;
    } else {
        if (x % 2 == 0)
            result = 2;     // positive even
        else
            result = 1;     // positive odd
    }
    return result;
}
\end{lstlisting}

\graphs{1}

{\small
\noindent\begin{tabularx}{\textwidth}{@{}l l l Y Y@{}}
\toprule
Block & idom & ipdom & Dominates & Post-dominates \\
\midrule
\texttt{entry} & -- & \texttt{if.end8} & all other blocks & -- \\
\texttt{if.then} & \texttt{entry} & \texttt{if.end8} & -- & -- \\
\texttt{if.else} & \texttt{entry} & \texttt{if.end7} & \texttt{if.then2, if.else3, if.then5, if.else6, if.end, if.end7} & -- \\
\texttt{if.then2} & \texttt{if.else} & \texttt{if.end7} & -- & -- \\
\texttt{if.else3} & \texttt{if.else} & \texttt{if.end} & \texttt{if.then5, if.else6, if.end} & -- \\
\texttt{if.then5} & \texttt{if.else3} & \texttt{if.end} & -- & -- \\
\texttt{if.else6} & \texttt{if.else3} & \texttt{if.end} & -- & -- \\
\texttt{if.end} & \texttt{if.else3} & \texttt{if.end7} & -- & \texttt{if.else3, if.then5, if.else6} \\
\texttt{if.end7} & \texttt{if.else} & \texttt{if.end8} & -- & \texttt{if.else, if.then2, if.else3, if.then5, if.else6, if.end} \\
\texttt{if.end8} & \texttt{entry} & \texttt{exit} & -- & all other blocks \\
\bottomrule
\end{tabularx}}

\medskip
\noindent\textbf{What stands out:} \texttt{if.end8} is the only way out,
so it post-dominates every block. But its immediate dominator is
\texttt{entry}, not \texttt{if.end7}, because the \texttt{x < 0} branch
jumps straight to it. The nested if-else is what gives both trees their
staircase shape.

%=====================================================================
\newpage
\section{Program 2: for loop with continue and break}

\begin{lstlisting}
int sumOdd(int n, int limit) {
    int sum = 0;
    for (int i = 1; i <= n; i++) {
        if (i % 2 == 0)
            continue;
        sum += i;
        if (sum > limit)
            break;
    }
    return sum;
}
\end{lstlisting}

\graphs{2}

{\small
\noindent\begin{tabularx}{\textwidth}{@{}l l l Y Y@{}}
\toprule
Block & idom & ipdom & Dominates & Post-dominates \\
\midrule
\texttt{entry} & -- & \texttt{for.cond} & all other blocks & -- \\
\texttt{for.cond} & \texttt{entry} & \texttt{for.end} & \texttt{for.body, if.then, if.end, if.then3, if.end4, for.inc, for.end} & \texttt{entry, if.then, if.end4, for.inc} \\
\texttt{for.body} & \texttt{for.cond} & \texttt{for.end} & \texttt{if.then, if.end, if.then3, if.end4, for.inc} & -- \\
\texttt{if.then} & \texttt{for.body} & \texttt{for.inc} & -- & -- \\
\texttt{if.end} & \texttt{for.body} & \texttt{for.end} & \texttt{if.then3, if.end4} & -- \\
\texttt{if.then3} & \texttt{if.end} & \texttt{for.end} & -- & -- \\
\texttt{if.end4} & \texttt{if.end} & \texttt{for.inc} & -- & -- \\
\texttt{for.inc} & \texttt{for.body} & \texttt{for.cond} & -- & \texttt{if.then, if.end4} \\
\texttt{for.end} & \texttt{for.cond} & \texttt{exit} & -- & all other blocks \\
\bottomrule
\end{tabularx}}

\medskip
\noindent\textbf{What stands out:} the \texttt{break} changes things.
Since \texttt{if.then3} jumps straight to \texttt{for.end}, the loop
header \texttt{for.cond} no longer post-dominates \texttt{for.body} or
\texttt{if.end}. The \texttt{continue} path (\texttt{if.then}) still goes
through \texttt{for.inc}, so \texttt{for.inc} post-dominates it.

%=====================================================================
\newpage
\section{Program 3: while loop with \texttt{\&\&} and an if-else}

\begin{lstlisting}
int collatzSteps(int n) {
    int steps = 0;
    while (n != 1 && steps < 1000) {
        if (n % 2 == 0)
            n = n / 2;
        else
            n = 3 * n + 1;
        steps++;
    }
    return steps;
}
\end{lstlisting}

\graphs{3}

{\small
\noindent\begin{tabularx}{\textwidth}{@{}l l l Y Y@{}}
\toprule
Block & idom & ipdom & Dominates & Post-dominates \\
\midrule
\texttt{entry} & -- & \texttt{while.cond} & all other blocks & -- \\
\texttt{while.cond} & \texttt{entry} & \texttt{land.end} & \texttt{land.rhs, land.end, while.body, if.then, if.else, if.end, while.end} & \texttt{entry, while.body, if.then, if.else, if.end} \\
\texttt{land.rhs} & \texttt{while.cond} & \texttt{land.end} & -- & -- \\
\texttt{land.end} & \texttt{while.cond} & \texttt{while.end} & \texttt{while.body, if.then, if.else, if.end, while.end} & \texttt{entry, while.cond, land.rhs, while.body, if.then, if.else, if.end} \\
\texttt{while.body} & \texttt{land.end} & \texttt{if.end} & \texttt{if.then, if.else, if.end} & -- \\
\texttt{if.then} & \texttt{while.body} & \texttt{if.end} & -- & -- \\
\texttt{if.else} & \texttt{while.body} & \texttt{if.end} & -- & -- \\
\texttt{if.end} & \texttt{while.body} & \texttt{while.cond} & -- & \texttt{while.body, if.then, if.else} \\
\texttt{while.end} & \texttt{land.end} & \texttt{exit} & -- & all other blocks \\
\bottomrule
\end{tabularx}}

\medskip
\noindent\textbf{What stands out:} the short-circuit \texttt{\&\&} splits
the loop test into three blocks. \texttt{land.rhs} can be skipped (when
\texttt{n == 1}), so it dominates nothing, while \texttt{land.end}, where
both halves meet, post-dominates \texttt{while.cond}. The if-else in the
body is the usual diamond: \texttt{while.body} dominates both arms and
\texttt{if.end} post-dominates them.

%=====================================================================
\newpage
\section{Program 4: do-while with a switch}

\begin{lstlisting}
int wordScore(const char *s) {
    int score = 0;
    int i = 0;
    do {
        switch (s[i]) {
        case 'a': case 'e': case 'i': case 'o': case 'u':
            score += 2;
            break;
        case ' ':
            break;
        default:
            score += 1;
        }
        i++;
    } while (s[i] != '\0');
    return score;
}
\end{lstlisting}

\graphs{4}

{\small
\noindent\begin{tabularx}{\textwidth}{@{}l l l Y Y@{}}
\toprule
Block & idom & ipdom & Dominates & Post-dominates \\
\midrule
\texttt{entry} & -- & \texttt{do.body} & all other blocks & -- \\
\texttt{do.body} & \texttt{entry} & \texttt{sw.epilog} & \texttt{sw.bb, sw.bb1, sw.default, sw.epilog, do.cond, do.end} & \texttt{entry} \\
\texttt{sw.bb} & \texttt{do.body} & \texttt{sw.epilog} & -- & -- \\
\texttt{sw.bb1} & \texttt{do.body} & \texttt{sw.epilog} & -- & -- \\
\texttt{sw.default} & \texttt{do.body} & \texttt{sw.epilog} & -- & -- \\
\texttt{sw.epilog} & \texttt{do.body} & \texttt{do.cond} & \texttt{do.cond, do.end} & \texttt{entry, do.body, sw.bb, sw.bb1, sw.default} \\
\texttt{do.cond} & \texttt{sw.epilog} & \texttt{do.end} & \texttt{do.end} & \texttt{entry, do.body, sw.bb, sw.bb1, sw.default, sw.epilog} \\
\texttt{do.end} & \texttt{do.cond} & \texttt{exit} & -- & all other blocks \\
\bottomrule
\end{tabularx}}

\medskip
\noindent\textbf{What stands out:} a do-while always runs its body once,
so \texttt{do.body}, \texttt{sw.epilog}, \texttt{do.cond} and
\texttt{do.end} form a straight chain in the dominator tree. All three
switch arms (the vowel case, the space case and default) meet at
\texttt{sw.epilog}, which makes it their immediate post-dominator.

%=====================================================================
\newpage
\section{Program 5: nested for loops with an early return}

\begin{lstlisting}
bool hasDuplicate(int a[], int n) {
    for (int i = 0; i < n; i++) {
        for (int j = i + 1; j < n; j++) {
            if (a[i] == a[j])
                return true;
        }
    }
    return false;
}
\end{lstlisting}

\graphs{5}

{\small
\noindent\begin{tabularx}{\textwidth}{@{}l l l Y Y@{}}
\toprule
Block & idom & ipdom & Dominates & Post-dominates \\
\midrule
\texttt{entry} & -- & \texttt{for.cond} & all other blocks & -- \\
\texttt{for.cond} & \texttt{entry} & \texttt{return} & \texttt{for.body, for.cond1, for.body3, if.then, if.end, for.inc, for.end, for.inc7, for.end9, return} & \texttt{entry, for.end, for.inc7} \\
\texttt{for.body} & \texttt{for.cond} & \texttt{for.cond1} & \texttt{for.cond1, for.body3, if.then, if.end, for.inc, for.end, for.inc7} & -- \\
\texttt{for.cond1} & \texttt{for.body} & \texttt{return} & \texttt{for.body3, if.then, if.end, for.inc, for.end, for.inc7} & \texttt{for.body, if.end, for.inc} \\
\texttt{for.body3} & \texttt{for.cond1} & \texttt{return} & \texttt{if.then, if.end, for.inc} & -- \\
\texttt{if.then} & \texttt{for.body3} & \texttt{return} & -- & -- \\
\texttt{if.end} & \texttt{for.body3} & \texttt{for.inc} & \texttt{for.inc} & -- \\
\texttt{for.inc} & \texttt{if.end} & \texttt{for.cond1} & -- & \texttt{if.end} \\
\texttt{for.end} & \texttt{for.cond1} & \texttt{for.inc7} & \texttt{for.inc7} & -- \\
\texttt{for.inc7} & \texttt{for.end} & \texttt{for.cond} & -- & \texttt{for.end} \\
\texttt{for.end9} & \texttt{for.cond} & \texttt{return} & -- & -- \\
\texttt{return} & \texttt{for.cond} & \texttt{exit} & -- & all other blocks \\
\bottomrule
\end{tabularx}}

\medskip
\noindent\textbf{What stands out:} the early \texttt{return true}.
Normally a loop's exit block post-dominates its header, but here
\texttt{for.end} does not post-dominate \texttt{for.cond1}, because you
can leave the inner loop through \texttt{if.then} and go straight to
\texttt{return}. So the ipdom of \texttt{for.cond1} is \texttt{return}
itself. Also, \texttt{return} is reached from both \texttt{if.then} and
\texttt{for.end9}, so its immediate dominator is the outer loop header
\texttt{for.cond}.

\end{document}
